\documentclass[11pt]{article}
\usepackage[margin=2.4cm]{geometry}
\usepackage{amsmath,amssymb}
\usepackage{microtype}
\pagestyle{empty}
\setlength{\parindent}{0pt}
\setlength{\parskip}{0.55em}

\begin{document}

\begin{center}
{\Large Your Figure~6, made solid}\\[2pt]
{\small the 2-Tamari decomposition of the associahedron
(\emph{Dyck}$^{(2)}_4$, from arXiv:1202.6269)}
\end{center}

\medskip

This model realizes the 14-cell decomposition of your hand-drawn figure with
every promise of the combinatorics made \emph{metric}: all 55 vertices of the
2-Tamari poset, all 14 cells convex with exactly planar walls, tiling a convex
3-dimensional associahedron exactly. The four combinatorial cubes are true
parallelepipeds, the eight pentagon cells true prisms (their pentagon pairs are
exact translates), so each cell is a genuine metric product of associahedra
$K_3$, $K_4$, $K_5$ --- the fibers of the projection
$\mathrm{Tam}^{(2)}_4 \to \mathrm{Tam}_4$. The $1$-skeleton is the Hasse
diagram of the 110 covering relations, each cell is exactly the interval of
its label, and a linear functional orients all 110 edges as the lattice order.

\textbf{The pieces come in repeated shapes.} The kit has only 10 distinct
shapes: four congruent pairs and six singletons. Two pairs are exact
translates. The other two pairs are congruent only by rotation --- and
provably so: a translation between them is \emph{impossible} in every valid
realization, because the covering orientations propagate through the
parallelogram walls with a fixed sign, and the correspondence the translation
would require reverses them.

\textbf{A theorem your figure taught us.} One naturally asks for this model in
the classical symmetric shape of the associahedron. It cannot be done:

\begin{quote}
\emph{No realization of the 14-cell decomposition with metric product cells
admits a symmetry of order greater than two. In particular the
$D_{3d}$-symmetric associahedron supports no such realization. The one
symmetry that can survive --- and is realized in the companion model --- is an
order-two rotation that inverts the Tamari order: the self-duality of the
Tamari lattice, as a motion of your hands.}
\end{quote}

The product cells force families of edges to be mutually parallel; a
$120^\circ$ symmetry would propagate these families until two edges sharing a
vertex became parallel, collapsing it. Turning the companion model by
$180^\circ$ returns the polytope to itself while every covering relation
reverses.

\begin{center}
\rule{0.42\textwidth}{0.4pt}\\[4pt]
{\small with admiration for a figure worth building}
\end{center}

\end{document}
